% BEDIENVORSCHAU MINIMALMODUS
% Öffentlicher Arbeitsablauf: eine Vorgabe, öffentliche Inhaltsbefehle und LetterBody.

\UseLetterForm{A}
\UsePreset{minimal-de}
\UseBackaddresstrue

\UseRecipientAddressbook{customer-service}
\SetLetterSubject{Beispiel für Form A im Modus: minimal}
\SetLetterOpening{Anrede}
\SetLetterClosing{Grußformel}

\begin{LetterBody}

\textbf{Was dieser Guide zeigt}

Wenn zusätzliche Informationsbereiche unnötig wären, hält \texttt{minimal-de} die Form-A-Seite bewusst ruhig. Das schwarze Monogramm bleibt sichtbar, während dekorative Farbblöcke und die automatische Anlagenliste ausgeblendet sind. Diese Vorschau aktiviert zusätzlich die optionale Rücksendezeile über dem Anschriftblock.

Lege Empfänger und Betreff in \texttt{content.tex} fest und ergänze die kurze Korrespondenz in \texttt{LetterBody}. Absenderinformationen kommen weiterhin aus \texttt{data/sender-data.tex}. Bleiben Demo-Daten unverändert, gibt das Toolkit vor der finalen Nutzung einen Hinweis im Compiler-Log aus, nicht als sichtbare Warnung im PDF. Die Vorgabe bietet einen an DIN 5008 Form A orientierten, nicht DIN-zertifizierten Ausgangspunkt und ist kein echtes Schreiben.

\end{LetterBody}
